% Options for packages loaded elsewhere
% Options for packages loaded elsewhere
\PassOptionsToPackage{unicode}{hyperref}
\PassOptionsToPackage{hyphens}{url}
\PassOptionsToPackage{dvipsnames,svgnames,x11names}{xcolor}
%
\documentclass[
]{article}
\usepackage{xcolor}
\usepackage[top=10pc,bottom=10pc,left=11pc,right=11pc,heightrounded,left=1.25in,right=1.25in]{geometry}
\usepackage{amsmath,amssymb}
\setcounter{secnumdepth}{-\maxdimen} % remove section numbering
\usepackage{iftex}
\ifPDFTeX
  \usepackage[T1]{fontenc}
  \usepackage[utf8]{inputenc}
  \usepackage{textcomp} % provide euro and other symbols
\else % if luatex or xetex
  \usepackage{unicode-math} % this also loads fontspec
  \defaultfontfeatures{Scale=MatchLowercase}
  \defaultfontfeatures[\rmfamily]{Ligatures=TeX,Scale=1}
\fi
\usepackage{lmodern}
\ifPDFTeX\else
  % xetex/luatex font selection
\fi
% Use upquote if available, for straight quotes in verbatim environments
\IfFileExists{upquote.sty}{\usepackage{upquote}}{}
\IfFileExists{microtype.sty}{% use microtype if available
  \usepackage[]{microtype}
  \UseMicrotypeSet[protrusion]{basicmath} % disable protrusion for tt fonts
}{}

\usepackage{color}
\usepackage{fancyvrb}
\newcommand{\VerbBar}{|}
\newcommand{\VERB}{\Verb[commandchars=\\\{\}]}
\DefineVerbatimEnvironment{Highlighting}{Verbatim}{commandchars=\\\{\}}
% Add ',fontsize=\small' for more characters per line
\usepackage{framed}
\definecolor{shadecolor}{RGB}{241,243,245}
\newenvironment{Shaded}{\begin{snugshade}}{\end{snugshade}}
\newcommand{\AlertTok}[1]{\textcolor[rgb]{0.68,0.00,0.00}{#1}}
\newcommand{\AnnotationTok}[1]{\textcolor[rgb]{0.37,0.37,0.37}{#1}}
\newcommand{\AttributeTok}[1]{\textcolor[rgb]{0.40,0.45,0.13}{#1}}
\newcommand{\BaseNTok}[1]{\textcolor[rgb]{0.68,0.00,0.00}{#1}}
\newcommand{\BuiltInTok}[1]{\textcolor[rgb]{0.00,0.23,0.31}{#1}}
\newcommand{\CharTok}[1]{\textcolor[rgb]{0.13,0.47,0.30}{#1}}
\newcommand{\CommentTok}[1]{\textcolor[rgb]{0.37,0.37,0.37}{#1}}
\newcommand{\CommentVarTok}[1]{\textcolor[rgb]{0.37,0.37,0.37}{\textit{#1}}}
\newcommand{\ConstantTok}[1]{\textcolor[rgb]{0.56,0.35,0.01}{#1}}
\newcommand{\ControlFlowTok}[1]{\textcolor[rgb]{0.00,0.23,0.31}{\textbf{#1}}}
\newcommand{\DataTypeTok}[1]{\textcolor[rgb]{0.68,0.00,0.00}{#1}}
\newcommand{\DecValTok}[1]{\textcolor[rgb]{0.68,0.00,0.00}{#1}}
\newcommand{\DocumentationTok}[1]{\textcolor[rgb]{0.37,0.37,0.37}{\textit{#1}}}
\newcommand{\ErrorTok}[1]{\textcolor[rgb]{0.68,0.00,0.00}{#1}}
\newcommand{\ExtensionTok}[1]{\textcolor[rgb]{0.00,0.23,0.31}{#1}}
\newcommand{\FloatTok}[1]{\textcolor[rgb]{0.68,0.00,0.00}{#1}}
\newcommand{\FunctionTok}[1]{\textcolor[rgb]{0.28,0.35,0.67}{#1}}
\newcommand{\ImportTok}[1]{\textcolor[rgb]{0.00,0.46,0.62}{#1}}
\newcommand{\InformationTok}[1]{\textcolor[rgb]{0.37,0.37,0.37}{#1}}
\newcommand{\KeywordTok}[1]{\textcolor[rgb]{0.00,0.23,0.31}{\textbf{#1}}}
\newcommand{\NormalTok}[1]{\textcolor[rgb]{0.00,0.23,0.31}{#1}}
\newcommand{\OperatorTok}[1]{\textcolor[rgb]{0.37,0.37,0.37}{#1}}
\newcommand{\OtherTok}[1]{\textcolor[rgb]{0.00,0.23,0.31}{#1}}
\newcommand{\PreprocessorTok}[1]{\textcolor[rgb]{0.68,0.00,0.00}{#1}}
\newcommand{\RegionMarkerTok}[1]{\textcolor[rgb]{0.00,0.23,0.31}{#1}}
\newcommand{\SpecialCharTok}[1]{\textcolor[rgb]{0.37,0.37,0.37}{#1}}
\newcommand{\SpecialStringTok}[1]{\textcolor[rgb]{0.13,0.47,0.30}{#1}}
\newcommand{\StringTok}[1]{\textcolor[rgb]{0.13,0.47,0.30}{#1}}
\newcommand{\VariableTok}[1]{\textcolor[rgb]{0.07,0.07,0.07}{#1}}
\newcommand{\VerbatimStringTok}[1]{\textcolor[rgb]{0.13,0.47,0.30}{#1}}
\newcommand{\WarningTok}[1]{\textcolor[rgb]{0.37,0.37,0.37}{\textit{#1}}}

\usepackage{longtable,booktabs,array}
\usepackage{calc} % for calculating minipage widths
% Correct order of tables after \paragraph or \subparagraph
\usepackage{etoolbox}
\makeatletter
\patchcmd\longtable{\par}{\if@noskipsec\mbox{}\fi\par}{}{}
\makeatother
% Allow footnotes in longtable head/foot
\IfFileExists{footnotehyper.sty}{\usepackage{footnotehyper}}{\usepackage{footnote}}
\makesavenoteenv{longtable}
\usepackage{graphicx}
\makeatletter
\newsavebox\pandoc@box
\newcommand*\pandocbounded[1]{% scales image to fit in text height/width
  \sbox\pandoc@box{#1}%
  \Gscale@div\@tempa{\textheight}{\dimexpr\ht\pandoc@box+\dp\pandoc@box\relax}%
  \Gscale@div\@tempb{\linewidth}{\wd\pandoc@box}%
  \ifdim\@tempb\p@<\@tempa\p@\let\@tempa\@tempb\fi% select the smaller of both
  \ifdim\@tempa\p@<\p@\scalebox{\@tempa}{\usebox\pandoc@box}%
  \else\usebox{\pandoc@box}%
  \fi%
}
% Set default figure placement to htbp
\def\fps@figure{htbp}
\makeatother


% definitions for citeproc citations
\NewDocumentCommand\citeproctext{}{}
\NewDocumentCommand\citeproc{mm}{%
  \begingroup\def\citeproctext{#2}\cite{#1}\endgroup}
\makeatletter
 % allow citations to break across lines
 \let\@cite@ofmt\@firstofone
 % avoid brackets around text for \cite:
 \def\@biblabel#1{}
 \def\@cite#1#2{{#1\if@tempswa , #2\fi}}
\makeatother
\newlength{\cslhangindent}
\setlength{\cslhangindent}{1.5em}
\newlength{\csllabelwidth}
\setlength{\csllabelwidth}{3em}
\newenvironment{CSLReferences}[2] % #1 hanging-indent, #2 entry-spacing
 {\begin{list}{}{%
  \setlength{\itemindent}{0pt}
  \setlength{\leftmargin}{0pt}
  \setlength{\parsep}{0pt}
  % turn on hanging indent if param 1 is 1
  \ifodd #1
   \setlength{\leftmargin}{\cslhangindent}
   \setlength{\itemindent}{-1\cslhangindent}
  \fi
  % set entry spacing
  \setlength{\itemsep}{#2\baselineskip}}}
 {\end{list}}
\usepackage{calc}
\newcommand{\CSLBlock}[1]{\hfill\break\parbox[t]{\linewidth}{\strut\ignorespaces#1\strut}}
\newcommand{\CSLLeftMargin}[1]{\parbox[t]{\csllabelwidth}{\strut#1\strut}}
\newcommand{\CSLRightInline}[1]{\parbox[t]{\linewidth - \csllabelwidth}{\strut#1\strut}}
\newcommand{\CSLIndent}[1]{\hspace{\cslhangindent}#1}



\setlength{\emergencystretch}{3em} % prevent overfull lines

\providecommand{\tightlist}{%
  \setlength{\itemsep}{0pt}\setlength{\parskip}{0pt}}





% -----------------------
% CUSTOM PREAMBLE STUFF
% -----------------------

% -----------------
% Typography tweaks
% -----------------
% Indent size
\setlength{\parindent}{1pc}  % 1p0

% Fix widows and orphans
\usepackage[all,defaultlines=2]{nowidow}

% List things
\usepackage{enumitem}
% Same document-level indentation for ordered and ordered lists
\setlist[1]{labelindent=\parindent}
\setlist[itemize]{leftmargin=*}
\setlist[enumerate]{leftmargin=*}

% Wrap definition list terms
% https://tex.stackexchange.com/a/9763/11851
\setlist[description]{style=unboxed}


% For better TOCs
\usepackage[titles]{tocloft}

% Remove left margin in lists inside longtables
% https://tex.stackexchange.com/a/378190/11851
\AtBeginEnvironment{longtable}{\setlist[itemize]{nosep, wide=0pt, leftmargin=*, before=\vspace*{-\baselineskip}, after=\vspace*{-\baselineskip}}}

% Allow for /singlespacing and /doublespacing
\usepackage{setspace}


% -----------------
% Title block stuff
% -----------------

% Abstract
\usepackage[overload]{textcase}
\usepackage[runin]{abstract}
\renewcommand{\abstractnamefont}{\sffamily\footnotesize\bfseries\MakeUppercase}
\renewcommand{\abstracttextfont}{\sffamily\small}
\setlength{\absleftindent}{\parindent * 2}
\setlength{\absrightindent}{\parindent * 2}
\abslabeldelim{\quad}
\setlength{\abstitleskip}{-\parindent}


% Keywords
\newenvironment{keywords}
{\vskip -3em \hspace{\parindent}\small\sffamily{\sffamily\footnotesize\bfseries\MakeUppercase{Keywords}}\quad}
{\vskip 3em}

  
% Title
\usepackage{titling}
\setlength{\droptitle}{3em}
\pretitle{\par\vskip 5em \begin{flushleft}\LARGE\sffamily\bfseries}
\posttitle{\par\end{flushleft}\vskip 0.75em}


% Authors
%
% PHEW this is complicated for a number of reasons!
%
% When using \and with multiple authors, the article class in LaTeX wraps each 
% author block in a tabluar environment with a hardcoded center alignment. It's 
% possible to use \preauthor{} to start tabulars with a left alignment {l}, but 
% that only applies to the first author because the others all use \and with the 
% hardcoded {c}. But we can override the \and command and add our own {l}
%
% (See https://github.com/rstudio/rmarkdown/issues/1716#issuecomment-560601691 
% for an example of redefining \and to just be \\)
%
% That's all great, except tabulars have some amount of default horizontal 
% padding, which makes left-aligned author blocks not actuall get fully 
% left-aligned on the page. We can set the horizontal padding for the column to 
% 0, but it requires some wonky syntax: {@{\hspace{0em}}l@{}}
\renewcommand{\and}{\end{tabular} \hskip 3em \begin{tabular}[t]{@{\hspace{0em}}l@{}}}
\preauthor{\begin{flushleft}
           \lineskip 1.5em 
           \begin{tabular}[t]{@{\hspace{0em}}l@{}}}
\postauthor{\end{tabular}\par\end{flushleft}}

% Omit the date since the \published command does that
\predate{}
\postdate{}

% Command for a note at the top of the first page describing the publication
% status of the paper.
\newcommand{\published}[1]{%
   \gdef\puB{#1}}
   \newcommand{\puB}{}
   \renewcommand{\maketitlehooka}{%
       \par\noindent\footnotesize\sffamily \puB}


% ------------------
% Section headings
% ------------------
\usepackage{titlesec}
\titleformat*{\section}{\Large\sffamily\bfseries\raggedright}
\titleformat*{\subsection}{\large\sffamily\bfseries\raggedright}
\titleformat*{\subsubsection}{\normalsize\sffamily\bfseries\raggedright}
\titleformat*{\paragraph}{\small\sffamily\bfseries\raggedright}

% \titlespacing{<command>}{<left>}{<before-sep>}{<after-sep>}
% Starred version removes indentation in following paragraph
\titlespacing*{\section}{0em}{2em}{0.1em}
\titlespacing*{\subsection}{0em}{1.25em}{0.1em}
\titlespacing*{\subsubsection}{0em}{0.75em}{0em}


% -----------
% Footnotes
% -----------
% NB: footmisc has to come after setspace and biblatex because of conflicts
\usepackage[bottom, flushmargin]{footmisc}
\renewcommand*{\footnotelayout}{\footnotesize}

\addtolength{\skip\footins}{10pt}    % vertical space between rule and main text
\setlength{\footnotesep}{5pt}  % vertical space between footnotes


% ----------
% Captions
% ----------
\usepackage[font={small,sf}, labelfont={small,sf,bf}]{caption}


% --------
% Macros
% --------
% pandoc will not convert text within \begin{} XXX \end{} to Markdown and will
% treat it as regular TeX. Because of this, it's impossible to do stuff like
% this:

% \begin{landscape}
%
% | One | Two   |
% |-----+-------|
% | my  | table |
% | is  | nice  |
%
% \end{landscape}
%
% Since it'll render like: | One | Two | |—–+——-| | my | table | | is | nice |
% 
% BUT, from this http://stackoverflow.com/a/41945462/120898 we can get around
% this by creating new commands for \begin and \end, like this:
\usepackage{pdflscape}
\newcommand{\blandscape}{\begin{landscape}}
\newcommand{\elandscape}{\end{landscape}}

% \blandscape
%
% | One | Two   |
% |-----+-------|
% | my  | table |
% | is  | nice  |
%
% \elandscape

% Same thing, but for generic groups
% But can't use \bgroup and \egroup because those are built-in TeX things
\newcommand{\stgroup}{\begingroup}
\newcommand{\fingroup}{\endgroup}


% ---------------------------
% END CUSTOM PREAMBLE STUFF
% ---------------------------
\makeatletter
\@ifpackageloaded{caption}{}{\usepackage{caption}}
\AtBeginDocument{%
\ifdefined\contentsname
  \renewcommand*\contentsname{Table of contents}
\else
  \newcommand\contentsname{Table of contents}
\fi
\ifdefined\listfigurename
  \renewcommand*\listfigurename{List of Figures}
\else
  \newcommand\listfigurename{List of Figures}
\fi
\ifdefined\listtablename
  \renewcommand*\listtablename{List of Tables}
\else
  \newcommand\listtablename{List of Tables}
\fi
\ifdefined\figurename
  \renewcommand*\figurename{Figure}
\else
  \newcommand\figurename{Figure}
\fi
\ifdefined\tablename
  \renewcommand*\tablename{Table}
\else
  \newcommand\tablename{Table}
\fi
}
\@ifpackageloaded{float}{}{\usepackage{float}}
\floatstyle{ruled}
\@ifundefined{c@chapter}{\newfloat{codelisting}{h}{lop}}{\newfloat{codelisting}{h}{lop}[chapter]}
\floatname{codelisting}{Listing}
\newcommand*\listoflistings{\listof{codelisting}{List of Listings}}
\makeatother
\makeatletter
\makeatother
\makeatletter
\@ifpackageloaded{caption}{}{\usepackage{caption}}
\@ifpackageloaded{subcaption}{}{\usepackage{subcaption}}
\makeatother
\usepackage{bookmark}
\IfFileExists{xurl.sty}{\usepackage{xurl}}{} % add URL line breaks if available
\urlstyle{same}
\hypersetup{
  pdftitle={Two R packages: one integrating pre-trained models for epidemiological time series forecasting, and another providing a curated epidemiological database to support efficient model reuse and development.},
  pdfauthor={Thiyanga S. Talagala},
  colorlinks=true,
  linkcolor={DarkSlateBlue},
  filecolor={Maroon},
  citecolor={DarkSlateBlue},
  urlcolor={DarkSlateBlue},
  pdfcreator={LaTeX via pandoc}}


% -----------------------
% END-OF-PREAMBLE STUFF
% -----------------------



% ---------------------- 
% Title block elements
% ---------------------- 
\usepackage{orcidlink}  % Create automatic ORCID icons/links

\title{Two R packages: one integrating pre-trained models for
epidemiological time series forecasting, and another providing a curated
epidemiological database to support efficient model reuse and
development.}


\author{
{\large Thiyanga S. Talagala~\orcidlink{0000-0002-0656-9789}}%
 \\%
Department of Statistics, Faculty of Applied Sciences, University of Sri
Jayewardenepura, Sri Lanka \\%
{\footnotesize \url{}} \and
}

\date{}


% Typeset URLs in the same font as their parent environment
%
% This has to come at the end of the preamble, after any biblatex stuff because 
% some biblatex styles (like APA) define their own \urlstyle{}
\usepackage{url}
\urlstyle{same}

% ---------------------------
% END END-OF-PREAMBLE STUFF
% ---------------------------
\begin{document}
% ---------------
% TITLE SECTION
% ---------------
\published{\textbf{2026-04-30}}

\maketitle



% -------------------
% END TITLE SECTION
% -------------------



\emph{This template was last revised on 2026-03-24.}

\section{Executive Summary}\label{executive-summary}

Total amount: USD 3500

Time series forecasting is essential for making informed decisions about
future. A recent trend in time series forecasting is pre-trained models.
Pre-trained models are machine learning models developed using massive
datasets, designed to be reused for similar tasks via transfer learning.
The two main advantages of pre-trained models are:

\begin{enumerate}
\def\labelenumi{\arabic{enumi}.}
\item
  Less data required: You do not need large datasets to train a model
  from scratch.
\item
  Faster training: The model has already trained and you need to tune
  the model for your purpose.
\end{enumerate}

This is very crucial when it is comes to time series forecasting as the
forecasts need to be given as quickly as possible. There are few
pre-trained models developed for time series forecasting. For example
\href{https://github.com/amazon-science/chronos-forecasting}{Chronos} by
Amazon, \href{https://github.com/google-research/timesfm}{TimesFM} by
Google Researchers, \href{https://github.com/Nixtla/nixtla}{TimeGPT} by
Nixtla,
\href{https://developer.ibm.com/tutorials/awb-foundation-model-time-series-forecasting/}{TinyTImeMixer}
by IBM.

In python there are several Python libraries supporting pre-trained
models. For example, \href{https://unit8co.github.io/darts/}{Darts},
\href{https://ts.gluon.ai/stable/}{GluonTS} and
\href{https://huggingface.co/models?other=time-series}{Hugging Face}.
However, there is no such package available in R that immediately pops
up that supports pre-trained models in time series forecasting.

There could be packages in R that use the idea of pre-trained model for
time series forecasting. For example
\href{https://cran.r-project.org/web/packages/seer/index.html}{seer} is
one such package. Here the authors use the phrase ``meta-learning''. One
disadvantage of this package is there is no way to tune the parameters
for the given time series. This package gives the functionalists to set
a pre-trained model not a pre-trained model.

Furthermore, most of the pre-trained models are trained based on
financial time series, sales time series, etc. \textbf{The purpose of
this project is to develop a simple pre-trained model and a database for
pre-training time series models, and to compile them into two R package
(One R package for the pre-trained model, another R package for the
pre-train model training database).} More specifically, in this research
we are planning to develop a pre-trained model for epidemiological
forecasting. This is useful in many different ways

\begin{enumerate}
\def\labelenumi{\arabic{enumi}.}
\item
  Give the idea of time series pre-trained model for R users. For
  example, in the field of time series forecasting one of the most
  referenced book is \href{https://otexts.com/fpp3/}{Forecasting
  Principles and Practice by Rob Hyndman and George Athanasopoulos}.
  This book has a both R version and a Python version. Link to the
  python varsion of the book is
  \href{https://otexts.com/fpppy/}{Forecasting: Principles and Practice,
  the Pythonic Way}. The python version of the book has a chapet called
  ``Foundation forecasting models''. This is not included in the R
  version. This highlights a gap in the filed. A light version a
  pre-trained model implemented in time series forecasting will help to
  fill the gap into some extent.
\item
  Most frontline health workers do not have expertise in time series
  forecasting; therefore, pre-trained models are essential to support
  their decision-making.
\item
  This provides a strong foundation to attract the attention of R users
  and developers toward pre-trained models.
\item
  In areas like computer vision and NLP, pre-trained models enable
  strong performance out-of-the-box. The lack of similar resources in
  R-based forecasting creates a gap in accessibility and efficiency.
\item
  Without pre-trained models, R users must design, tune, and validate
  models from scratch, which requires advanced statistical and
  computational knowledge.
\end{enumerate}

R is evolving toward more accessible, automated, and scalable workflows,
but pre-trained models and plug-and-play forecasting systems are still
limited, especially in applied domains like epidemiology.

\section{Signatories}\label{signatories}

\subsection{Project team}\label{project-team}

\href{https://thiyanga.netlify.app/}{Thiyanga S. Talagala}

\begin{itemize}
\item
  Data collection
\item
  Data cleaning
\item
  Pre-trained model development
\item
  Documentation
\end{itemize}

\subsection{Contributors}\label{contributors}

None

\subsection{Consulted}\label{consulted}

None

\section{The Problem}\label{the-problem}

\subsection{What the problem is}\label{what-the-problem-is}

There is currently no R package that provides ready-to-use pre-trained
model for time series forecasting.

\subsection{Who it affects}\label{who-it-affects}

There is a rich tidyverse(\citeproc{ref-tidyverse}{Wickham et al. 2019})
friendly time series analysis pipeline called tidyverts. This contains
tsibble (\citeproc{ref-tsibble}{Wang, Cook, and Hyndman 2020}), fable
(\citeproc{ref-fable}{O'Hara-Wild, Hyndman, and Wang 2024a}), fabletools
(\citeproc{ref-fabletools}{O'Hara-Wild, Hyndman, and Wang 2024b}) R
packages that helps for data structures for storing time series data,
functions to do data wrangling, and functions for time series
forecasting. Pre-trained models are widely used across many domains
today. This is a very effective approach for time series forecasting
particularly with limited data. Even though R has already established a
rich ecosystem for time series forecasting, the unavailability of
pre-trained models makes R users to shift to alternative programming
languages or limit their ability to use this efficient method for time
series forecasting.

\subsection{Why is it a problem}\label{why-is-it-a-problem}

It becomes a problem because the absence of pre-trained models changes
forecasting from a quick, accessible task into a slow, expertise-heavy
process. This also makes R users underutilization of advanced methods.
This forces R users to stick to forecasts models such as ARIMA, ETS,
Neural Networks, Prophet, etc. In otherowrds, this limits the
accessibility of modern forecasting methods for R users.

\subsection{What will solving the problem enable (why should it be
solved)}\label{what-will-solving-the-problem-enable-why-should-it-be-solved}

In this work we will develop a light weight pre-trained model for
epidemiology time series forecasting. To train the pre-trained model I
will use epidemiology counts data published by health ministries,
institutions, etc. across multiple countries. For example, Epidemiology
Unit, Ministry of Health, Sri Lanka published weekly epidemilogy reports
as of
\href{https://www.epid.gov.lk/weekly-epidemiological-report/weekly-epidemiological-report}{here}
corresponds to the following diseases:

\begin{itemize}
\item
  Dengue fever
\item
  Dysentery
\item
  Encephalitis
\item
  Encephalitis fever
\item
  Food Poison
\item
  Leptospirosis
\item
  Typhus
\item
  Viral Hepatitis
\item
  H. Rabies
\item
  Chickenpox
\item
  Meningitis
\item
  Leishman
\item
  Tubeculosis
\item
  Leprosy
\end{itemize}

I will be using all these weekly counts to my pre-trained model
database. Inaddition to I will collect data from the following websites:

\href{https://vizhub.healthdata.org/gbd-results/}{Institute for Health
Metrics and Evaluation}

\href{https://www.who.int/teams/global-tuberculosis-programme/data}{WHO
Tuberculosis database}

\url{https://opendengue.org/}

I will also explore freely available open-access databases for data
collection. Crrently, there is no single global website that contains
clean ``ready-made counts'' for all these diseases together.

\subsection{Brief summary of existing work and previous attempts (e.g.,
relevant R
packages)}\label{brief-summary-of-existing-work-and-previous-attempts-e.g.-relevant-r-packages}

Previous packages:

\texttt{denguedatahub} R package: I compiled dengue data for 216
countries and compiled into a package called denguedatahub. The project
website is at \url{https://denguedatahub.netlify.app/}. This project is
funded by RConortium ISC 2022 grant. Since 2022, I am maintaining this
project. This dataset will also be used train the odel.

\texttt{DISC}: The R package I developed
\url{https://github.com/SMART-Research/DISC} contains disese counts for
Dysentery, Encephalitis, Encephalitis fever, Food Poison, Leptospirosis,
Typhus, Viral Hepatitis, H. Rabies, Chickenpox, Meningitis, Leishman,
Tubeculosis and Leprosy. This will also be used in pre-trained model
development.

\texttt{seer} package: The seer package provides implementations of a
novel framework for forecast model selection using time series features.
We call this framework FFORMS (Feature-based FORecast Model Selection).
For more details see our
\href{https://onlinelibrary.wiley.com/doi/abs/10.1002/for.2963}{paper}.
The R package repository is: \url{https://github.com/thiyangt/seer}.
This is a joint work with Prof Rob Hyndman and Prof George
Athanasopoulos

\texttt{vegcast}: Feature engineering to train a pre-trained model for
vegetable price forecasting. This work is part of the project funded by
\url{https://forecasters.org/programs/research-awards/forecasting-for-social-good-research-grant/}.
I am the 2025 F4SG grant receipient as listed on the website. The github
repositoy of the project \url{https://github.com/thiyangt/vegcast}. The
project repository is
\url{https://thiyangt.github.io/pricecast.vegetables/}.

\texttt{tsdataleaks}: R package to identify data leaks in time series
data available at: \url{https://github.com/thiyangt/tsdataleaks}

\texttt{vegetablesSriLanka}: R package for vegetable prices in Sri
Lanka. Link \url{https://github.com/thiyangt/vegetablesSriLanka}

In addition to there are many other packages I developed and published
on CRAN.

\subsection{If proposing changes to R itself: letter of support from R
Core
member}\label{if-proposing-changes-to-r-itself-letter-of-support-from-r-core-member}

No changes to R

\section{The proposal}\label{the-proposal}

\subsection{Overview}\label{overview}

There is no single global website that contains clean ``ready-made
counts'' for all these diseases together. Furthermore, there is no R
package to access pre-trained model in time series forecasting. In this
project we develop a pre-trained model for epidemiology time series
forecasting.

\subsection{Detail}\label{detail}

The project can be devided into three phases.

\textbf{Phase 1:} Preparation of database to develop pre-trained model

Based on the database an R package will be developed named
\texttt{EpiDataHub} -- Curated Epidemiological Time Series Repository

\textbf{Phase 2:} Train a pre-trained model

\textbf{Phase 3:} Development of two R packages corresponding to work of
Phase 1 and Phase 2.

\textbf{Pre-trained modelling steps can be listed as follows:}

Step 1: Train a global model

\begin{Shaded}
\begin{Highlighting}[]
\NormalTok{global\_data }\OtherTok{\textless{}{-}} \FunctionTok{read.csv}\NormalTok{(}\StringTok{"all\_districts\_data.csv"}\NormalTok{)}

\NormalTok{rf\_global }\OtherTok{\textless{}{-}} \FunctionTok{randomForest}\NormalTok{(}
\NormalTok{  disease\_cases }\SpecialCharTok{\textasciitilde{}}\NormalTok{ lag\_cases }\SpecialCharTok{+}\NormalTok{ disease\_type }\SpecialCharTok{+}\NormalTok{ year }\SpecialCharTok{+}\NormalTok{ month }\SpecialCharTok{+}\NormalTok{ week }\SpecialCharTok{+}\NormalTok{ day }\SpecialCharTok{+} 
\NormalTok{  climate\_type }\SpecialCharTok{+}\NormalTok{ area,}
  \AttributeTok{data =}\NormalTok{ global\_data,}
  \AttributeTok{ntree =} \DecValTok{1000}
\NormalTok{)}
\end{Highlighting}
\end{Shaded}

Step 2: Save model

\begin{Shaded}
\begin{Highlighting}[]
\FunctionTok{saveRDS}\NormalTok{(rf\_global, }\StringTok{"rf\_global\_model.rds"}\NormalTok{) }
\end{Highlighting}
\end{Shaded}

This model will be converted into an R package called
\texttt{EpiCastNet} -- Pre-trained Forecasting Models for Epidemiology.

Step 3: Tune the model for a new given time series to generate forecasts
for the new series.

First generate forecast for the new series using global model. That is
our pre-trained model

\begin{Shaded}
\begin{Highlighting}[]
\NormalTok{local\_data}\SpecialCharTok{$}\NormalTok{global\_pred }\OtherTok{\textless{}{-}} \FunctionTok{predict}\NormalTok{(rf\_global, }\AttributeTok{newdata =}\NormalTok{ local\_data)}
\end{Highlighting}
\end{Shaded}

Develop separate local model to the new series and use the previous
prediction from the global model as a predictor. This step is refereed
to as tuning the model.

\begin{Shaded}
\begin{Highlighting}[]
\NormalTok{rf\_local }\OtherTok{\textless{}{-}} \FunctionTok{randomForest}\NormalTok{(}
\NormalTok{  dengue\_cases }\SpecialCharTok{\textasciitilde{}}\NormalTok{ rainfall }\SpecialCharTok{+}\NormalTok{ temperature }\SpecialCharTok{+}\NormalTok{ humidity }\SpecialCharTok{+}
\NormalTok{    lag\_cases }\SpecialCharTok{+}\NormalTok{ global\_pred,}
  \AttributeTok{data =}\NormalTok{ local\_data,}
  \AttributeTok{ntree =} \DecValTok{500}
\NormalTok{)}
\end{Highlighting}
\end{Shaded}

The final forecast is a combination from global model, local model. Here
the concept of staking is used to tune the model. Stacking (or stacked
learning / stacked generalization) is a method where the prediction from
one model is used as an input (feature) for another model.

In summary

Model 1: Pre-trained model

Learns general/global patterns

↓

Its prediction becomes a new variable

↓

Model 2: Tune pre-trained model

Uses that prediction + original predictors to make a better final
prediction

\subsubsection{Minimum Viable Product}\label{minimum-viable-product}

Development of two R packages:

\begin{enumerate}
\def\labelenumi{\roman{enumi}.}
\item
  A database of epidemiological time series to support the training of
  pre-trained models. The R package name is \texttt{EpiDataHub} --
  Curated Epidemiological Time Series Repository
\item
  A pre-trained model package for time series forecasting. The R package
  name is \texttt{EpiCastNet} -- Pre-trained Forecasting Models for
  Epidemiology.
\end{enumerate}

\subsubsection{Architecture}\label{architecture}

\subsubsection{Assumptions}\label{assumptions}

Sufficient high-quality epidemiological time series data will be
available from open-access sources to build a reliable database.

If sufficient data are not available, feature-based data augmentation
methods will be used to simulate data
(\citeproc{ref-kang2020gratis}{Kang, Hyndman, and Li 2020}).

\subsubsection{External dependencies}\label{external-dependencies}

The project relies on several established R packages for modeling and
time series analysis, including tidyverse, randomForest, tidymodels,
xgboost, tsibble, fabletools, and fable. These are the main packages
that will be used in the development.

\section{Project plan}\label{project-plan}

\subsection{Start-up phase}\label{start-up-phase}

\subsection{Technical delivery}\label{technical-delivery}

The work on the package development is scheduled for May 2026 - April
2027.

\begin{longtable}[]{@{}
  >{\raggedright\arraybackslash}p{(\linewidth - 2\tabcolsep) * \real{0.2857}}
  >{\raggedright\arraybackslash}p{(\linewidth - 2\tabcolsep) * \real{0.7143}}@{}}
\toprule\noalign{}
\begin{minipage}[b]{\linewidth}\raggedright
Timeline
\end{minipage} & \begin{minipage}[b]{\linewidth}\raggedright
Tasks and milestones
\end{minipage} \\
\midrule\noalign{}
\endhead
\bottomrule\noalign{}
\endlastfoot
May, June, July 2026 & Development of database to train pre-trained
model \\
August, 2026 & Pre-train model development and cross-validate \\
September, October 2026 & Development of the package \\
November 2026 & Mid Progress Review and Mid project presentation to get
feedback \\
December 2026 & Continue the development work based on the feedback \\
January 2027 & Write and share a blog post about the work, release the
first version of the package \\
February 2027 & Write a manual for the package and update pre-trained
model using new data \\
March 2027 & Released second version of the package. Presentation to
health care officers \\
End of April 2027 & Final presentation \\
\end{longtable}

\subsection{Other aspects}\label{other-aspects}

A separate project website will be created to publish blog post, to
updae about the news, such as release of the first version of the
pre-trained model, etc.

In addition, LinkedIn, X, BlueSky social media platforms will be used to
increase the visibility of the work.

An R Ladies Colombo event will be organized to promote work.

I am also planning to present the work at R Medicine conference by R
Consortium.

\subsection{Budget \& funding plan}\label{budget-funding-plan}

Total budget required for the project: USD 3500

\begin{enumerate}
\def\labelenumi{\arabic{enumi}.}
\item
  Data collection: USD 2000

  Collect data from different organizations magazines, websites, etc.

  Write codes for web scrapping to extract data from websites .

  Data quality analysis.

  Write codes to clean data and convert them into analysis ready tidy
  format.
\item
  Development of the pre-trained model: USD 1000

  Train pre-trained model

  Hyper parameter-tuning

  Time-series cross validation
\item
  Write documentations: USD 500

  Project website development

  Write blog post

  Write package vignette (User manual)

  Talk/Workshop for health care/ decision-makers officers on how to use
  pre-trained model for epidemiology forecasting
\end{enumerate}

\section{Success}\label{success}

\subsection{Definition of done}\label{definition-of-done}

The project's definition of done includes the following deliverables:

\begin{enumerate}
\def\labelenumi{\arabic{enumi}.}
\item
  Development of a project website, with regularly published blog posts
  and news updates.
\item
  Development of two R packages:

  \begin{enumerate}
  \def\labelenumii{\roman{enumii}.}
  \item
    A database of epidemiological time series to support the training of
    pre-trained models. The R package name is \texttt{EpiDataHub} --
    Curated Epidemiological Time Series Repository
  \item
    A pre-trained model package for time series forecasting. The R
    package name is \texttt{EpiCastNet} -- Pre-trained Forecasting
    Models for Epidemiology.
  \end{enumerate}
\end{enumerate}

Both packages should be published on CRAN.

\begin{enumerate}
\def\labelenumi{\arabic{enumi}.}
\setcounter{enumi}{2}
\tightlist
\item
  Completion and delivery of the final project presentation.
\end{enumerate}

\subsection{Measuring success}\label{measuring-success}

The project will be considered a success if the following criteria are
met:

\begin{itemize}
\tightlist
\item
  The key deliverables are met: the package includes the planned
  functionality, which is documented and tested.
\item
  The documentation is available on the package website.
\item
  At least 75\% of the open issues on the planned work are closed or at
  least become work-in-progress.
\item
  Feedback is sought (and incorporated as far as possible) from the R
  community.
\end{itemize}

\subsection{Future work}\label{future-work}

Future work will focus on continuously updating the epidemiological time
series database and annually releasing updated versions of the
pre-trained forecasting models. Additional efforts will include
expanding the database to cover more diseases and regions, improving
model accuracy through ongoing fine-tuning, and incorporating user
feedback to enhance functionality and usability of the packages.

Furthermore, by presenting this work at conferences, I plan to find
strong collaborators for the project.

\phantomsection\label{refs}
\begin{CSLReferences}{1}{0}
\bibitem[\citeproctext]{ref-kang2020gratis}
Kang, Yanfei, Rob J Hyndman, and Feng Li. 2020. {``GRATIS: GeneRAting
TIme Series with Diverse and Controllable Characteristics.''}
\emph{Statistical Analysis and Data Mining: The ASA Data Science
Journal} 13 (4): 354--76.

\bibitem[\citeproctext]{ref-fable}
O'Hara-Wild, Mitchell, Rob Hyndman, and Earo Wang. 2024a. \emph{Fable:
Forecasting Models for Tidy Time Series}.
\url{https://doi.org/10.32614/CRAN.package.fable}.

\bibitem[\citeproctext]{ref-fabletools}
---------. 2024b. \emph{Fabletools: Core Tools for Packages in the
'Fable' Framework}.
\url{https://doi.org/10.32614/CRAN.package.fabletools}.

\bibitem[\citeproctext]{ref-tsibble}
Wang, Earo, Dianne Cook, and Rob J Hyndman. 2020. {``A New Tidy Data
Structure to Support Exploration and Modeling of Temporal Data.''}
\emph{Journal of Computational and Graphical Statistics} 29 (3):
466--78. \url{https://doi.org/10.1080/10618600.2019.1695624}.

\bibitem[\citeproctext]{ref-tidyverse}
Wickham, Hadley, Mara Averick, Jennifer Bryan, Winston Chang, Lucy
D'Agostino McGowan, Romain François, Garrett Grolemund, et al. 2019.
{``Welcome to the {tidyverse}.''} \emph{Journal of Open Source Software}
4 (43): 1686. \url{https://doi.org/10.21105/joss.01686}.

\end{CSLReferences}




\end{document}
